\chapter{Технологический раздел}

\section{Средства реализации}
Алгоритмы для данной лабораторной работы были реализованы на языке C++, при использовании компилятора gcc версии 13.3.1, так как в стандартной библиотеке приведенного языка
есть функция \texttt{clock\_gettime}, которая (при использовании макропеременной \texttt{CLOCK\_THREAD\_CPUTIME\_ID}) позволяет рассчитать процессорное время конкретного потока \cite{cpp-time}.

\section{Реализация алгоритмов}
Листинги исходных кодов программ  \ref{lst:lev_alg.cpp}--\ref{lst:dam_lev_alg.cpp} приведены в приложении.

\section{Тестирование}
Для написания модульных тестов использовалась библиотека $gtest$ \cite{gtests}. Данные тесты приведены в листинге \ref{lst:tests.cpp}.

Тесты рассмотрены в таблице \ref{t:unit_tests}.

\begin{table}[!ht]
	\begin{center}
		\small
		\begin{threeparttable}
			\caption{Модульные тесты}
			\label{t:unit_tests}
			\begin{tabular}{|c|c|c|c|c|c|}
				\hline
				\multicolumn{2}{|c|}{\bfseries Входные данные}
				& \multicolumn{4}{c|}{\bfseries Расстояние и алгоритм} \\ 
				\hline 
				&
				& \multicolumn{1}{c|}{\bfseries Дамерау---Левенштейна} 
				& \multicolumn{3}{c|}{\bfseries Левенштейна} \\ \cline{3-6}
				
				\bfseries S1 & \bfseries S2 & \bfseries Линейный & \bfseries Линейный
				
				& \multicolumn{2}{c|}{\bfseries Рекурсивный} \\ \cline{5-6}
				& & & & \bfseries Без кеша & \bfseries С кешем \\
				\hline
				hello & hello & 0 & 0 & 0 & 0 \\
				\hline
				meoww & hello & 4 & 4 & 4 & 4 \\
				\hline
				ppppp & hello & 5 & 5 & 5 & 5 \\
				\hline
				helol & hello & 1 & 2 & 2 & 2 \\
				\hline
			\end{tabular}	
		\end{threeparttable}
	\end{center}
\end{table}

Все тесты были успешно пройдены.
